% Part 2 导言：几何、材质与动态
\vspace{1em}
\begin{center}
\textcolor{techgray}{\small 本部分将3DGS从``好看''推向``准确''和``动态''——几何表面重建、材质分解与逆渲染、动态场景与4D建模。}
\end{center}
\vspace{0.5em}

\noindent\textbf{本部分包含：}
\begin{itemize}[leftmargin=2em,itemsep=1pt]
\item[\textcolor{techgreen}{\textbullet}] \textbf{第4章}（几何重建）— 深度/法线正则化、SDF引导、网格提取、2D Surfel
\item[\textcolor{techgreen}{\textbullet}] \textbf{第5章}（材质与光照）— 反射建模、材质分解、重光照、HDR、PBR
\item[\textcolor{techgreen}{\textbullet}] \textbf{第6章}（动态与4D）— 变形场、4D高斯、单目动态重建、运动建模、场景分解
\end{itemize}

\noindent\textbf{阅读建议：}第4章和第5章相对独立，可按需选读。第6章是后续多个应用章节的基础——自动驾驶（第18章）和机器人（第17章）的动态建模技术都源于本章。

\noindent\textbf{与其他部分的联系：}第4章的SDF方法与第3章的基元改进互补；第5章的材质分解结果可直接用于第9章的场景编辑；第6章的变形场技术是第7章物理仿真的表示基础。

\vspace{1em}
